\documentclass[11pt,a4paper]{article}
\usepackage{amsmath,amssymb,amsthm,mathtools,booktabs,hyperref}
\usepackage[margin=27mm]{geometry}
\setlength{\emergencystretch}{1em}
\newtheorem{theorem}{Theorem}[section]
\newtheorem{lemma}[theorem]{Lemma}
\newtheorem{corollary}[theorem]{Corollary}
\newtheorem{proposition}[theorem]{Proposition}
\newcommand{\Q}{\mathbb Q}\newcommand{\Z}{\mathbb Z}\newcommand{\N}{\mathbb N}
\title{An unrestricted-degree undecidability transfer in $9m+2$ rational unknowns}
\author{\(\Z\Q\) Diophantine Project}
\date{23 November 2026}
\begin{document}\maketitle
\begin{abstract}
Suppose an integer polynomial in $m$ total rational coordinates, including its
input, has input projection exactly $\Z$. We prove that rational solvability of
polynomial equations is then undecidable in $9m+2$ unknowns. The construction
combines the published nine-unknown mixed-domain theorem of Sun with a ternary
quadratic expression whose integral values over the rationals are exactly the
nonnegative integers. Substituting this expression directly removes its output
coordinate. For the project's current defining-polynomial premises this gives
227 unknowns in the general route, 245 in the fully specified route, and 101
under ordinary rational abc. These are unrestricted-degree conclusions. Their
arithmetic assumptions are precisely those of the chosen definition of the
integers; no unconditional claim about those inherited assumptions is made.
\end{abstract}

\section{What is counted}
Write $\N=\{0,1,2,\ldots\}$. A defining polynomial is a fixed
$F\in\Z[T,Y_1,\ldots,Y_{m-1}]$ such that, for every $t\in\Q$,
\begin{equation}\label{defining}
 t\in\Z\quad\Longleftrightarrow\quad
 \exists\mathbf y\in\Q^{m-1}\ F(t,\mathbf y)=0.
\end{equation}
Thus $m$ includes the input coordinate $T$. If the existence or defining property
of $F$ is proved only under specified premises, all consequences below inherit
those premises. The transfer arguments themselves introduce no number-theoretic
conjecture.

For undecidability, the input to the decision problem is the finite coefficient
list of a polynomial, and the unknowns are its rational arguments. A natural
number $a$ coding a decision instance is substituted into the coefficients
before asking for a rational solution; it is not a quantified unknown. If one
instead retains $a$ as a displayed variable of the universal family, one must
add one to the counts in this paper.

\section{A ternary wrapper with rational witnesses}
Define
\[
 L(s,t,u)=s^2+t^2+u^2+u.
\]
\begin{lemma}\label{wrapper}
For every $n\in\Z$,
\[
 n\ge0\quad\Longleftrightarrow\quad
 \exists s,t,u\in\Q\quad n=L(s,t,u).
\]
Moreover, witnesses can be chosen integral in the forward direction.
\end{lemma}
\begin{proof}
If $n\ge0$, then $4n+1$ is not of the form $4^a(8b+7)$. The
Gauss--Legendre three-square theorem therefore gives
$4n+1=A^2+B^2+C^2$ with $A,B,C\in\Z$. Exactly one of the three
squares is odd, as is seen by reducing modulo four. Reorder them so that
$A=2s$, $B=2t$, and $C=2u+1$. Subtracting one and dividing by four gives
the required identity with integral witnesses. Conversely, for rational
$s,t,u$ one has
\[
 L(s,t,u)=s^2+t^2+(u+\tfrac12)^2-\tfrac14\ge-\tfrac14.
\]
If this value is an integer, it is nonnegative.
\end{proof}

The requirement that $L(s,t,u)$ be integral is indispensable; its rational
values need not be nonnegative. Formula~\eqref{defining} supplies exactly that
requirement in the construction below. The witnesses $s,t,u$ themselves
need no integrality certificates.

\begin{proposition}[Optimality among quadratic wrappers]\label{wrapperoptimal}
There is no rational polynomial $P\in\Q[S,T]$ of degree at most two such that
\[
 P(\Q^2)\cap\Z=\N.
\]
Consequently the three rational coordinates of Lemma~\ref{wrapper} cannot be
replaced by two while retaining a single quadratic substitution and the same
integral-value criterion.
\end{proposition}
\begin{proof}
Suppose such a $P$ exists. Since zero belongs to its image, choose
$\mathbf a\in\Q^2$ with $P(\mathbf a)=0$. After multiplying the argument by a
sufficiently divisible positive integer $M$, all coefficients of
$P(\mathbf a+M\mathbf X)$ are integral. Thus its values at $\Z^2$ are integers.
If the quadratic homogeneous part has a negative direction, it has a rational
negative direction by continuity, and hence an integral one after scaling.
The values along that integral line are negative for sufficiently large
arguments, a contradiction. The quadratic part is therefore positive
semidefinite. Its kernel is a rational subspace. A nonzero linear part on
that kernel would again give arbitrarily negative integral values along an
integral line. Hence the linear part vanishes on the kernel.

It follows by solving a rational linear system and completing the square
that $P=c+q$, after a rational affine change of coordinates, where $c\in\Q$
and $q$ is a positive semidefinite rational quadratic form. If $q$ has rank
zero the image is constant, so it cannot contain $\N$. If its rank is one,
write $q=aU^2$ with $a>0$ rational. At any prime $p$ outside a finite set,
every nonzero rational value of $q$ has even $p$-adic valuation. Choose a
nonnegative integer $n$ with $v_p(n-c)=1$, which is possible by prescribing
one residue class modulo $p^2$. Then $n$ is not represented.

If $q$ has rank two, rational diagonalization writes
$q=a(U^2+dV^2)$ with $a,d>0$ rational. The quadratic field
$\Q(\sqrt{-d})$ has infinitely many inert primes, by the usual quadratic
case of Dirichlet's theorem on primes in arithmetic progressions. Choose
one outside the numerator and denominator primes of $a,d$ and the
denominator primes of $c$, and different
from two. At this prime $-d$ is a nonsquare. For any nonzero rational pair
$(U,V)$, put $e=\min\{v_p(U),v_p(V)\}$. Dividing by $p^{2e}$ leaves a unit,
since a nontrivial zero of $U^2+dV^2$ modulo $p$ is impossible. Consequently
$v_p(q)=2e$ is even. Again choose $n\in\N$ with $v_p(n-c)=1$. This
contradicts representation of every nonnegative integer.
\end{proof}

\section{The mixed-domain universal input}
We use the following published theorem, identifying the imported mathematical
input rather than presenting it as a new proof of the DPRM theorem.

\begin{theorem}[Sun's mixed-domain theorem]\label{sun}
For every recursively enumerable $A\subseteq\N$ there is an integer polynomial
$Q_A(a,z_1,\ldots,z_8,n)$ such that, for every $a\in\N$,
\[
 a\in A\quad\Longleftrightarrow\quad
 \exists z_1,\ldots,z_8\in\Z\ \exists n\in\N\quad
 Q_A(a,z_1,\ldots,z_8,n)=0.
\]
\end{theorem}
This is a direct consequence of \cite[Theorem 1.1(i), equations (1.2)--(1.3)]{Sun}.
Sun proves the stronger forward statement that all nine integer witnesses can
be required to exceed any prescribed bound. A later construction and detailed
formalization appear in \cite[Theorem 3 and Definition 5.3]{BDHMS}. The theorem
has eight unrestricted integer witnesses and one nonnegative integer witness;
treating all nine as positive unnecessarily increases the cost over $\Q$.

\section{The transfer theorem and its proof}
\begin{theorem}\label{main}
Assume \eqref{defining}. For every recursively enumerable $A\subseteq\N$ there is
an integer-polynomial family $G_A(a,\mathbf v)$ with at most $9m+2$ rational
unknowns $\mathbf v$ such that, for all $a\in\N$,
\[
 a\in A\quad\Longleftrightarrow\quad
 \exists\mathbf v\in\Q^{9m+2}\ G_A(a,\mathbf v)=0.
\]
In particular, no algorithm decides rational solvability for every integer
polynomial in at most $9m+2$ unknowns.
\end{theorem}
\begin{proof}
Choose $Q_A$ as in Theorem~\ref{sun}. Introduce eight rational coordinates
$z_1,\ldots,z_8$, three rational coordinates $s,t,u$, and nine independent
blocks $\mathbf y^{(i)}\in\Q^{m-1}$, $1\le i\le9$. Put
\begin{align}\label{construction}
 G_A={}&Q_A(a,z_1,\ldots,z_8,L(s,t,u))^2\notag\\
 &+\sum_{i=1}^{8}F(z_i,\mathbf y^{(i)})^2
   +F(L(s,t,u),\mathbf y^{(9)})^2.
\end{align}
This is an integer polynomial. Its number of rational coordinates is
\[
 8+3+9(m-1)=9m+2.
\]
Suppose first that $G_A=0$ at a rational point. Every displayed summand
is a square in the ordered field $\Q$, so every squared expression is zero.
By \eqref{defining}, $z_1,\ldots,z_8$ and $L(s,t,u)$ are integers.
Lemma~\ref{wrapper} gives $L(s,t,u)\ge0$. The remaining equation is precisely
a witness to $a\in A$ in Theorem~\ref{sun}.

Conversely, suppose $a\in A$. Choose the integer witnesses $z_1,\ldots,z_8,n$
from Theorem~\ref{sun}. By Lemma~\ref{wrapper}, choose integral $s,t,u$ with
$n=L(s,t,u)$. Apply \eqref{defining} independently to the nine integers
$z_1,\ldots,z_8,n$, obtaining all nine rational witness blocks. Every term
in \eqref{construction} is zero, so $G_A$ has a rational zero.

Finally take $A$ recursively enumerable but not recursive. Once $F$ and $Q_A$
are fixed, the map $a\mapsto G_A(a,\mathbf v)$ is computable, because it is
substitution into a fixed polynomial with integer coefficients. A decision
algorithm for rational solvability in the displayed number of unknowns
would decide membership in $A$, a contradiction.
\end{proof}

The fixed polynomial $F$ need not have an effectively known construction
parameter for this existence form of undecidability. Every fixed finite
integer coefficient list is computable and can be built into a Turing machine.
An explicitly evaluable reduction requires an explicitly specified defining
member; an ineffective existence assertion alone must not be described as an
implemented numerical reduction.

\begin{proposition}[General mixed-domain count]\label{mixed}
Suppose a set is represented by a polynomial over $r$ unrestricted integer
unknowns and $s$ nonnegative integer unknowns. Under \eqref{defining}, it has a
rational-polynomial representation with
\[
 rm+s(m+2)=(r+s)m+2s
\]
rational unknowns. The same count holds with strictly positive integer unknowns
by replacing $L$ with $1+L$.
\end{proposition}
\begin{proof}
For each unrestricted integer coordinate retain its input and use one
$(m-1)$-coordinate block for $F$. For each nonnegative coordinate substitute
an independent copy of $L$ directly into the source polynomial and into
its copy of $F$. This costs three rational coordinates and $m-1$ witnesses.
The equivalence follows exactly as in Theorem~\ref{main}; for strict positivity,
integer values of $1+L$ are at least one and every positive integer occurs.
\end{proof}

\section{Degree bounds and source qualifications}
If $F$ has degree $D$ and the mixed-domain source polynomial has degree $d$,
then the elementary bound obtained from \eqref{construction} is
\[
 \deg G_A\le4\max(D,d).
\]
A better recorded bound uses the explicit source of \cite{BDHMS}. There the
polynomial obtained after substituting $n=L(s,t,u)$ already has degree at most
\begin{equation}\label{eta}
 \eta(\nu,\delta)=15616+233856\delta
   +233952\delta(2\delta+1)^{\nu+1}
   +467712\delta^2(2\delta+1)^{\nu+1},
\end{equation}
whenever $(\nu,\delta)$ is a universal pair over the nonnegative integers.
This is the same substitution used in their Section 6.1. Therefore
\begin{equation}\label{degree}
 \deg G_A\le\max\{2\eta(\nu,\delta),4D\}.
\end{equation}
Using the published $(58,4)$ pair gives
\begin{align*}
 \eta(58,4)
 &=1681043235226619916301182624511918527834137733707408448335539840,\\
 2\eta(58,4)
 &=3362086470453239832602365249023837055668275467414816896671079680.
\end{align*}
The exact integer has been recomputed directly from \eqref{eta}.

For completeness, \cite{BDHMS} also states the smaller numerical value
\[
 \eta(32,12)=950817549694171759711025515571236610412597656252821888.
\]
The authors expressly note that Jones mentions $(32,12)$ without giving its
proof. We do not promote this smaller degree to the principal audited bound.
This qualification does not affect the variable count in Theorem~\ref{main}.
The present note does not claim to have rerun the authors' Isabelle development.
The formulas and theorem statements were checked in version 3 of their paper,
dated 22 October 2026; its degree formula agrees with version 2.

\section{Application to the current project routes}
\begin{center}
\begin{tabular}{lrrl}\toprule
Defining route & $m$ & Rational unknowns & Status inherited\\\midrule
General few-variable &25&227&Retained arithmetic premises\\
Fully specified general &27&245&Retained arithmetic premises\\
Ordinary rational abc &11&101&abc and retained premises\\\bottomrule
\end{tabular}
\end{center}
The input counts and degree bounds used here are those of the project's
v88 defining-polynomial ledger: respectively $D\le48318382080$,
$D\le183368679424$, and $D\le3261569$. In every case $4D<2\eta(58,4)$,
so \eqref{degree} gives one common explicit degree upper bound above.
These are bounds for undecidability, not lower bounds or claims of optimality.
They are not quartic results. Replacing the large-degree expressions by
quadratic circuits adds coordinates and requires a separate count.

The ordinary-abc route is deliberately retained as a separate row. The
no-additional-abc routes still inherit the arithmetic interfaces of their
chosen $F$, so the numerical conclusions must not be called unconditional
merely because abc is absent. An independent proof of those interfaces
would immediately make the corresponding numerical application unconditional.

\section{A precise obstruction to a tempting compression}
The following elementary observation rules out replacing arbitrary tuples
of integer inputs by fewer homogeneous integral-valued tests. It does not
rule out systems with additional constraints or shared witness constructions.

\begin{proposition}\label{homogeneous}
Let $1\le k<r$ and let
$H_1,\ldots,H_k\in\Z[X_1,\ldots,X_r]$ be homogeneous polynomials of positive
degrees. There are infinitely many primes $p$ and tuples
$\mathbf q\in\Z[1/p]^r\setminus\Z^r$ such that
$H_j(\mathbf q)\in\Z$ for every $j$.
\end{proposition}
\begin{proof}
Over the algebraic closure of $\Q$, fewer than $r$ positive-degree homogeneous
equations in $r$ variables have a common nonzero zero. One way to see this
is the projective dimension theorem; equivalently, the ideal generated by
$k$ elements has height at most $k$, whereas the irrelevant maximal ideal
has height $r$. Choose such a zero $\boldsymbol\alpha$ and normalize one
nonzero coordinate to one. Its coordinates lie in a number field $K$.

There are infinitely many rational primes for which $K$ embeds into $\Q_p$.
For a short justification choose an integral primitive element of $K$ with
minimal polynomial $P\in\Z[T]$. Infinitely many primes divide nonzero values
of $P$ (Schur's elementary theorem). To recall the elementary argument for Schur's assertion, if all these prime
divisors belonged to a finite set $S$, choose $a$ with $A=P(a)\ne0$ and put
$Q(T)=P(a+AT)/A\in\Z[T]$. Then $Q(0)=1$. For a multiple $M$ of all
primes in $S$, every $Q(Mk)$ is congruent to one at those primes, whereas
$|Q(Mk)|>1$ for some $k$. A prime divisor of that value is outside $S$,
a contradiction. Except for finitely many primes dividing
the discriminant, each such divisor gives a simple root of $P$ modulo $p$,
which lifts by Hensel's lemma to a root in $\Q_p$ and hence an embedding
$K\hookrightarrow\Q_p$. Exclude also the finitely many denominator primes
of the coordinates of $\boldsymbol\alpha$ in an integral basis. All its
coordinates then lie in $\Z_p$, and the normalized coordinate is still one.

Let $D=\max_j\deg H_j$. Choose an integer vector $\mathbf a$ congruent to
the embedded vector $\boldsymbol\alpha$ modulo $p^D$, choosing the normalized
coordinate to be exactly one. Then $H_j(\mathbf a)$ is divisible by $p^D$.
For $d_j=\deg H_j\le D$, homogeneity gives
\[
 H_j(\mathbf a/p)=p^{-d_j}H_j(\mathbf a)\in\Z.
\]
The normalized coordinate of $\mathbf a/p$ is $1/p$, so this tuple is not
integral. All but finitely many primes supplied above work.
\end{proof}

In particular, an integral value of a sum of even powers or of a homogeneous
norm form cannot, without other constraints, force all of two or more rational
arguments to be integers. This explains why replacing the fifteen input
certificates in the earlier quartic construction by one such expression
is not a valid shortcut. The proposition does not assert a lower bound for
arbitrary definitions or arbitrary universal systems.


\section{Why the ninth integer certificate cannot simply be omitted}
The mixed-domain proof uses the final integer witness to enforce divisibility,
in addition to positivity and square conditions. Keeping a rational ternary
wrapper while discarding its integrality test loses that information.
Here is an exact counterexample to that attempted use of the relation-combining
lemma. The polynomial $M_3$ used in \cite{Sun,BDHMS} is the product over all
sign choices of the expressions
\[
 \begin{gathered}
 S^2n+T^2-S^2(2R-1)
 \left(T^2+X^3+\varepsilon_1\sqrt{A_1}
       +\varepsilon_2X\sqrt{A_2}+\varepsilon_3X^2\sqrt{A_3}\right),
 \\
 X=1+A_1^2+A_2^2+A_3^2.
 \end{gathered}
\]
Take $A_1=A_2=A_3=1$, $R=1$, $S=3$, and $T=1$. Then $X=4$,
and the all-positive factor vanishes at
\[
 n=1+64+1+4+16-\frac19=\frac{773}{9}
   =L\left(\frac{17}{3},\frac{22}{3},0\right).
\]
Nevertheless $S\nmid T$. Thus the sign condition supplied by $L$ and
rational solvability of $M_3=0$ do not suffice for the integral
relation-combining theorem. This example is a counterexample to the proposed
general lemma; it does not claim to realize all the additional source-circuit
relations simultaneously.

\begin{thebibliography}{9}
\bibitem{Sun} Z.-W. Sun, \emph{Further results on Hilbert's Tenth Problem},
Science China Mathematics \textbf{64} (2021), 281--306.
\url{https://arxiv.org/abs/1704.03504}; version 7, 28 January 2021.
\bibitem{BDHMS} J. Bayer, M. David, M. Ha\ss ler, Y. Matiyasevich, and D. Schleicher,
\emph{Diophantine equations over $\mathbb Z$: universal bounds and parallel formalization},
\url{https://arxiv.org/abs/2506.20909v3}, 22 October 2026.
Theorem 3, Definition 5.3, Theorem A, and Section 6.1.
\bibitem{Jones} J. P. Jones, \emph{Universal Diophantine equation},
Journal of Symbolic Logic \textbf{47} (1982), 549--571,
\url{https://doi.org/10.2307/2273588}.
\end{thebibliography}
\end{document}
